% =====================================================================
%  Appendix A -- the source document's own section map.
%  Source: tier2.md:113-191 (the Contents table and the collapsed
%  "Full section list"), reproduced so that the source's own
%  Part 0 / I / II / III / IV / V / VI / VII architecture survives the
%  re-architecture into a physics arc.
% =====================================================================

\section{The source document's section map}
\label{app:sectionmap}

This report re-architects \code{study/tier2.md} into a physics arc
(\hyperref[sec:status]{Status of this document}). The source's own organisation --- Part~0 the astrophysics,
then node by node from S7 to S10, then policy, register and method --- is
reproduced here so that nothing about it is lost, and so that a reader who knows
the source document can navigate this one.
Appendix~\ref{app:concordance} is the heading-by-heading map between them.

\subsection{The source's Contents table}
\label{app:sourcecontents}

\begingroup
\small
\begin{longtable}{@{}L{2.2cm} L{12.8cm}@{}}
\toprule
\textbf{Part} & \textbf{Covers} \\
\midrule
\endfirsthead
\toprule
\textbf{Part} & \textbf{Covers} \\
\midrule
\endhead
\bottomrule
\endfoot

\textbf{0} --- Astrophysical prerequisites &
The equilibrium hierarchy (kinetics $\to$ QSE $\to$ NSE) and its history;
silicon burning as photodisintegration rearrangement; the \Ye{} $\to$ observable
chain derived end to end \textbf{and the derivative nobody has assembled}; why
conservation is structural; \textbf{what consumes the output and what speedup is
achievable (the $1/f$ ceiling)}; the kill-test as a matrix question \\

\textbf{I} (S7) --- Stoichiometry \& conservation &
$\nu$ as a linear map; conservation laws \textbf{are} the left null space (with
the rank measurement proving $C$ is complete); \textbf{the same object in CRN
theory --- deficiency, weak reversibility, the flux cone}; lepton ledgers and
the tautology trap; Target A's theorem \textbf{and what it does not guarantee
(positivity)}; the projector derived twice; the \textbf{Target-A $\equiv$
Target-B reachability theorem}; \textbf{the effective dimension of the reachable
set and where \Ye{} hides}; the graph and $K$ \\

\textbf{II} (S8) --- Fluxes \& cancellation &
The pair map and the net identity; \kap{} four ways --- DB diagnostic, condition
number, \textbf{thermodynamic affinity}, \textbf{entropy production};
\textbf{the entropy-weighted active set, the second-law constraint on a learned
flux, and the Lyapunov bound}; the weak $\kap \equiv 1$ contract; the store; the
handshake \textbf{and the $\Delta t/\tau$ conditioning variable it implies}; the
vacuous pass \\

\textbf{III} (S9) --- NSE and QSE &
Statistical equilibrium from the grand canonical ensemble;
$\mu_i = Z\mu_p + N\mu_n$ derived; the Saha coefficient term-by-term; the Newton
system and its analytic Jacobian; \textbf{the monotonicity proof}; QSE clusters
and \uG; QSE as a slow manifold; the Guidry $\delta$ criterion; the empty-mask
verdict \\

\textbf{IV} (S10) --- Stiff integration &
Stiffness from the Jacobian spectrum; BDF and L-stability; the augmented
$[Y, \Phifl]$ state; \textbf{why \Phifl{} cannot be recovered --- and why that
makes the flux head unidentifiable under a $\Delta Y$-only loss}; the analytic
$\partial R/\partial Y$; the energy identity as a \emph{data} check \textbf{and
the Wegscheider measurement that explains its residual}; the cost wall \\

\textbf{V} --- Policy \& measurement discipline &
The structural invariants encoded in code that refuses --- then the three the
repo does \textbf{not} encode: \textbf{the sampling measure on the regime box},
\textbf{cluster-robust inference (657k samples $\approx$ 300 independent ones)},
and verification practice (convergence order, manufactured solutions, trivial
baselines) \\

\textbf{VI} --- The open register &
\textbf{The single register: 25 items (R-01 \ldots R-25), each pointing at the
section that derives it, plus a ranked top eight.} R-24 (learning theory) and
R-25 (shelf life) are written out here, having no home section \\

\textbf{VII} --- How the register was built &
The two gap-audit passes and the axes each swept; \textbf{thirteen candidate
axes checked and found already covered}; is the audit converging? (yes, with
three qualifications) \\
\end{longtable}
\endgroup

\subsection{The source's full section list}
\label{app:sourcelist}

Reproduced verbatim from the source's collapsed ``Full section list'' block.

\begin{srclistingtiny}
0.1  The equilibrium hierarchy          NSE / QSE / kinetics as flux statements
     counting DOF . the historical arc . what the project bets on it
0.2  Silicon burning, mechanically      why ~3 GK . two clusters, one bridge set
                                        the Fe-peak endpoint . not a ladder
0.3  The Ye chain, end to end           M_ch ~ Ye^2 . deleptonisation . the observable
     0.3.4 the error budget -> the 3e-6 gate
     0.3.5 (!) the OUTPUT end: dM_Ni/dYe has never been assembled
0.4  Why conservation is structural      left null vectors . the soft-penalty failure
0.5  What breaks on a leak
0.5b Deployment                          what consumes the output . imposed vs fed-back
                                         (!) the cost ceiling is 1/f, not the raw speed
0.6  The decision this tier serves        0.7 what is NOT here . 0.8 self-check

I.1  nu as a linear map                  I.8  The projector derived twice
I.2  Conservation = left null space      I.9  Extended space & the laundering trap
     the rank measurement                I.10 Target A = Target B (reachability)
     why float64 is exact here           I.10b (!) effective dimension; Ye in the tail
I.2b CRN theory, deficiency, flux cone   I.11 The bipartite graph, radius, K
     deficiency 126/327 - the DZT does   I.12 Code walkthrough
     NOT apply, measured, and why        I.13 Oracles: the 30-test gate
I.3-I.6 weak sector . ledgers . b+ . C   I.14 self-check
I.7  Target A's theorem
I.7b (!) ...and what it does NOT guarantee: positivity

II.1 From lambda to R to f+/f-           II.6  The store: lossless derivation
II.2 The pair map                        II.7  The handshake: premise & tolerance
II.3 The net identity, proved            II.7b (!) condition on dt/tau, not log dt
II.4 kappa four ways: DB . condition number  II.8  Cancellation-aware residuals
     . affinity . entropy production        II.9  Measured: the vacuous pass
II.4b (!) sigma = 2k_B s kappa artanh kappa II.10 Code + oracles
      -> entropy-weighted active set        II.11 self-check
      -> sign(phi-hat) from the affinity
      -> the Lyapunov bound in F (S20)      II.5 weak columns: kappa = 1 structurally
      ALL strong/EM-only: sigma, A and F
      are undefined on weak columns

III.1-III.7  Saha derived; the two constraints; Newton + analytic Jacobian;
             (!) the dYe/du_p > 0 proof; batching & the scalar fallback
III.8-III.11 QSE clusters & u_G; the slow manifold; delta vs r_QSE; the group boundary
III.12-III.14 the three negative results; code + oracles; self-check

IV.1 Stiffness from the spectrum         IV.6  Conservation by construction
IV.2 Why BDF, and what L-stability buys  IV.7  Error control for Phi
IV.3 The augmented state [Y, Phi]        IV.8  The energy identity is a DATA check
IV.4 (!) Why Phi is not recoverable      IV.8b (!) Wegscheider: Q _|_ null(nu) FAILS
IV.4b (!) ...so the flux head is               rms 0.55 keV, max 3.5 keV
      unidentifiable in 528 directions         predicts the stock-MESA kappa floor
IV.5 dR/dY by the product rule           IV.9  The cost wall, diagnosed
                                         IV.10 Code + oracles . IV.11 self-check

V.1  Invariants encoded in code that refuses
V.2  (!) the sampling MEASURE on the box (not its support): 35/62/82% within .1/.2/.3 dex
V.3  (!) cluster-robust inference: ICC 0.22-0.70, deff 6.5-18.6, n_eff ~ 300
V.4  V&V: convergence order . manufactured solutions . NSE-as-a-baseline

VI.1 The register, R-01 ... R-25         VII.1 The method: two passes, which axes
VI.2 R-24 learning theory                VII.2 13 axes checked, 13 already covered
     R-25 shelf life                     VII.3 Is the audit converging?
VI.3 Ranked: the top eight
\end{srclistingtiny}

\begin{aside}
\textbf{Glyph substitutions in the listing above.} The block is set in a
monospaced face, so a handful of source glyphs were transliterated rather than
risk a missing glyph: $\nu \to$ \code{nu}, $\kap \to$ \code{kappa},
$\sigma \to$ \code{sigma}, $\aff \to$ \code{A}, $\hat\netflux \to$
\code{phi-hat}, $\Phifl \to$ \code{Phi}, $\Delta t \to$ \code{dt},
$\tau \to$ \code{tau}, $\delta \to$ \code{delta}, $\Ye \to$ \code{Ye},
$\beta^+ \to$ \code{b+}, $\Mch \propto \Ye^2 \to$ \code{M\_ch \textasciitilde{}
Ye\textasciicircum2}, $\equiv\ \to$ \code{=}, $\to\ \to$ \code{->},
$\perp\ \to$ \code{\_|\_}, the interpunct $\to$ \code{.}, and the warning glyph
to \code{(!)}. Nothing else was changed: the wording, the line breaks, the
two-column layout, and the ordering are the source's --- only the column
padding of the \S II.4\,/\,\S II.4b stanza was widened, because transliterating
$\kap$ and $\sigma$ to their names lengthens the left column.
\end{aside}
